% =============================================================================
% Rapport technique — Deepfake Adaptive Fusion (MLOps)
% Compilation : ./scripts/build_report.sh   (ou pdflatex + bibtex + pdflatex x2)
% =============================================================================
\documentclass[11pt,a4paper]{article}

\usepackage[utf8]{inputenc}
\usepackage[T1]{fontenc}
\usepackage[french]{babel}
\usepackage{lmodern}
\usepackage[margin=2.4cm]{geometry}
\usepackage{graphicx}
\usepackage{booktabs}
\usepackage{amsmath,amssymb}
\usepackage{float}
\usepackage{caption}
\usepackage{subcaption}
\usepackage{xcolor}
\usepackage{hyperref}
\usepackage{listings}
\usepackage{array}
\usepackage{enumitem}
\usepackage{siunitx}

\graphicspath{{../reports/figures/}{../docs/}{./generated/}}
\hypersetup{colorlinks=true, linkcolor=blue!50!black, urlcolor=blue!50!black,
            citecolor=blue!50!black, pdftitle={Deepfake Adaptive Fusion — Rapport technique},
            pdfauthor={Projet MLOps}}
\lstset{basicstyle=\ttfamily\small, breaklines=true, frame=single,
        backgroundcolor=\color{gray!5}, columns=fullflexible}

\title{\textbf{Deepfake Adaptive Fusion}\\[0.3em]
\large Détection de deepfakes par fusion adaptative spatiale et fréquentielle\\
\large Industrialisation MLOps : DVC, MLflow, ONNX, FastAPI, React, Docker, CI/CD}
\author{Projet MLOps \& NLP}
\date{\today}

\begin{document}
\maketitle

\begin{abstract}
Ce rapport décrit la conception, l'implémentation et l'évaluation d'un détecteur
de deepfakes combinant un flux spatial convolutif et un flux fréquentiel, agrégés
par une porte adaptative. Le prototype issu d'un notebook Colab a été refactorisé
en un dépôt reproductible : pipeline DVC, suivi MLflow, export ONNX quantifié,
API FastAPI, interface React interactive, images Docker et CI/CD GitHub Actions.
Les métriques présentées sont générées par les artefacts du dépôt
(\texttt{reports/}, MLflow) ; aucune valeur n'est codée en dur. Nous discutons
aussi les limites de généralisation et les risques d'usage.
\end{abstract}

\tableofcontents
\newpage

% -----------------------------------------------------------------------------
\section{Contexte et problématique}
Les deepfakes, images ou vidéos synthétiques générées par réseaux antagonistes
génératifs (GAN) ou modèles de diffusion, posent des risques de désinformation,
de fraude et d'atteinte à la réputation. La détection est un problème
d'apprentissage supervisé binaire : distinguer une image réelle d'une image
générée. Les approches spatiales capturent des artefacts locaux (yeux, textures,
bords), tandis que les approches fréquentielles exploitent les traces spectrales
laissées par le sur-échantillonnage et les opérations de génération
\cite{wang2020cnn,durall2020watch,frank2020leveraging}.

La problématique industrielle est double : (i) construire un modèle performant et
interprétable ; (ii) le rendre reproductible, testé, déployable et maintenable.
Le projet répond aux deux, avec une attention particulière à l'anti-fuite, à la
traçabilité des expériences et à l'honnêteté des conclusions.

% -----------------------------------------------------------------------------
\section{Données, EDA et anti-fuite}
\subsection{Dataset principal}
Le dataset \texttt{xhlulu/140k-real-and-fake-faces} \cite{kaggle140k} fournit des
visages réels et générés (StyleGAN) avec des splits officiels
\textit{train/valid/test}. Pour tenir sur un budget Colab, un sous-échantillonnage
déterministe par classe est appliqué (seed 42), paramétré dans
\texttt{params.yaml}. Les tailles complètes restent accessibles en mettant les
limites à \texttt{null}.

\subsection{EDA}
L'analyse exploratoire vérifie l'équilibre des classes, la résolution, les
formats et le spectre de Fourier moyen réel \textit{vs} faux. Les figures sont
produites dans \texttt{reports/figures/} (spectre moyen, distribution des scores,
matrice de confusion, courbes ROC). Aucune figure n'est fabriquée pour le
rapport : les blocs conditionnels ci-dessous n'affichent que les artefacts
réellement générés.

\subsection{Anti-fuite}
Trois protections sont implémentées :
\begin{enumerate}[nosep]
  \item splits officiels, sans mélange manuel ;
  \item sous-échantillonnage par classe avec graine fixe ;
  \item test MD5 des contenus entre splits (\texttt{tests/test\_no\_leakage.py}),
        et découpage par vidéo/identité pour les jeux externes.
\end{enumerate}
Les augmentations ne sont appliquées qu'au train, après le split. FFHQ est exclu
des réels externes car déjà présent dans 140k.

% -----------------------------------------------------------------------------
\section{Architecture et justification du flux fréquentiel}
Le modèle \texttt{AdaptiveFusionDetector} comprend :
\begin{itemize}[nosep]
  \item \texttt{spatial\_net} : backbone convolutif (ConvNeXt-T \cite{liu2022convnext})
        sur l'image RGB ;
  \item \texttt{freq\_net} : CNN sur \texttt{freq\_input = log1p(|fftshift(fft2(x))|)} ;
  \item \texttt{spatial\_projection} et \texttt{frequency\_projection} : projections
        vers un espace commun ;
  \item \texttt{fusion\_gate} : porte softmax qui produit
        $(w_{\text{spatial}}, w_{\text{fréquentiel}})$ ;
  \item \texttt{classifier} : logit binaire ; \texttt{residual\_scale} : échelle
        résiduelle apprise.
\end{itemize}
La fusion est adaptative : le modèle apprend, image par image, la confiance à
accorder à chaque modalité. Les variantes \texttt{spatial\_only},
\texttt{freq\_only} et \texttt{concat\_no\_gate} servent de baselines et
d'ablations.

\paragraph{Export ONNX.} \texttt{torch.fft.fft2} ne s'exporte pas proprement.
Nous adoptons l'Option A : la FFT est calculée hors du graphe, côté inférence
(NumPy), et le modèle ONNX prend deux entrées (\texttt{image}, \texttt{freq\_input})
et renvoie trois sorties (\texttt{logit}, \texttt{w\_spatial},
\texttt{w\_frequency}). La parité PyTorch/ONNX est testée
($\max|\Delta \text{logit}| < 10^{-3}$, 100\,\% de concordance des classes).
Voir \texttt{docs/DECISIONS.md}.

% -----------------------------------------------------------------------------
\section{Protocole expérimental}
\paragraph{Recette.} AdamW avec learning rates différenciés (backbone, fusion,
classifieur), schedule cosinus, label smoothing, gradient clipping, AMP,
early stopping sur le F1 de validation, reprise via \texttt{last\_checkpoint.pt}.

\paragraph{Runs obligatoires.} À données, seed et schedule identiques :
R1 \texttt{spatial\_only}, R2 \texttt{freq\_only}, R3 \texttt{concat\_no\_gate},
R4 \texttt{fusion}. Chaque run est tracé dans MLflow (params, métriques par
époque, artefacts, tags git SHA/auteur/axe).

\paragraph{Ablations.} RandAugment \cite{cubuk2020randaugment} + MixUp/CutMix
\cite{zhang2018mixup,yun2019cutmix} (train uniquement), Focal Loss
\cite{lin2017focal} \textit{vs} BCE + label smoothing, gel/dégel progressif
\textit{vs} LR différenciés, robustesse JPEG/flou/redimensionnement.

\paragraph{Cross-dataset.} Les jeux externes (Celeb-DF v2 \cite{li2020celebdf},
FF-C23 \cite{rossler2019faceforensics}) ne servent qu'au test, jamais à
l'entraînement. Le protocole est identique (même prétraitement, même seuil) et
les résultats sont rapportés séparément.

\paragraph{Seuil.} Seuil par défaut $P(\text{fake}) \ge 0{,}5$, calibré sur la
validation ; la courbe précision/rappel et le seuil retenu sont écrits dans
\texttt{reports/threshold\_calibration.json} et exposés par l'API.

% -----------------------------------------------------------------------------
\section{Résultats}
\subsection{Métriques in-distribution}
Les valeurs ci-dessous sont générées par
\texttt{python -m src.models.evaluate} à partir de \texttt{reports/metrics.json}.

\IfFileExists{generated/metrics_table.tex}
  {\input{generated/metrics_table.tex}}
  {\begin{quote}\itshape Tableau des métriques non généré : lancer
   \texttt{python -m src.models.evaluate} puis
   \texttt{python scripts/export\_report\_tables.py}.\end{quote}}

\subsection{Baseline vs fusion et ablations}
\IfFileExists{generated/comparison_table.tex}
  {\input{generated/comparison_table.tex}}
  {\begin{quote}\itshape Tableau comparatif non généré : lancer
   \texttt{python scripts/run\_experiments.py}.\end{quote}}

\paragraph{Conclusion honnête.} La fusion n'apporte un gain que si les tableaux
ci-dessus le montrent, \emph{en particulier en cross-dataset}. Si le gain
in-distribution ne se transfère pas, nous le documentons sans le masquer. Le
modèle final est choisi sur la validation, jamais sur le test.

\subsection{Cross-dataset et robustesse}
Les résultats cross-dataset (\texttt{reports/cross\_dataset.json}) et la
dégradation sous JPEG/flou/redimensionnement (\texttt{reports/robustness.json})
sont générés par l'évaluation. Ils quantifient la généralisation réelle, souvent
plus faible qu'in-distribution.

\IfFileExists{../reports/figures/confusion_matrix_test.png}{%
\begin{figure}[H]\centering
\includegraphics[width=0.55\linewidth]{confusion_matrix_test.png}
\caption{Matrice de confusion sur le test (artefact réel).}
\end{figure}}{}

\IfFileExists{../reports/figures/roc_curve_test.png}{%
\begin{figure}[H]\centering
\includegraphics[width=0.55\linewidth]{roc_curve_test.png}
\caption{Courbe ROC sur le test (artefact réel).}
\end{figure}}{}

% -----------------------------------------------------------------------------
\section{Interprétabilité}
\paragraph{Grad-CAM.} La carte Grad-CAM \cite{selvaraju2017gradcam} est calculée
sur la dernière feature map du flux spatial, pour la classe cible
(\texttt{fake} ou \texttt{real}), puis superposée à l'image avec une opacité
réglable. Le backend utilise PyTorch en chargement paresseux pour
\texttt{/predict/visualize}, ONNX Runtime restant utilisé pour
\texttt{/predict}.

\paragraph{Poids de fusion.} $w_{\text{spatial}}$ et $w_{\text{fréquentiel}}$
sont exposés dans chaque prédiction. Ils permettent de voir si le modèle
s'appuie plutôt sur la texture spatiale ou sur les traces spectrales ; une
analyse par classe est produite dans \texttt{reports/}.

\IfFileExists{../reports/figures/gradcam_examples.png}{%
\begin{figure}[H]\centering
\includegraphics[width=0.85\linewidth]{gradcam_examples.png}
\caption{Exemples Grad-CAM (artefact réel).}
\end{figure}}{}

% -----------------------------------------------------------------------------
\section{Optimisation ONNX (FP32 / FP16 / INT8)}
Trois variantes sont produites : FP32, FP16 (conversion de poids) et INT8
(quantification dynamique \texttt{onnxruntime.quantization}). Pour chacune :
taille du fichier, latence CPU batch 1 (médiane et p95 après warm-up),
accuracy/F1/AUC sur le test. Le compromis est résumé dans le tableau généré par
l'évaluation.

\paragraph{Règle de décision.} Si INT8 dégrade l'AUC de plus d'un point, il n'est
pas déployé ; FP16/FP32 est conservé et le compromis est documenté.

\IfFileExists{../reports/figures/onnx_benchmark.png}{%
\begin{figure}[H]\centering
\includegraphics[width=0.7\linewidth]{onnx_benchmark.png}
\caption{Latence et taille des variantes ONNX (artefact réel).}
\end{figure}}{}

% -----------------------------------------------------------------------------
\section{Architecture logicielle et MLOps}
\paragraph{Dépôt.} Le code est séparé en \texttt{src/} (données, modèles,
utilitaires), \texttt{backend/} (FastAPI), \texttt{frontend/} (React/Vite),
\texttt{scripts/}, \texttt{tests/} et \texttt{docs/}. Le notebook nettoyé
importe \texttt{src/} au lieu de dupliquer le code.

\paragraph{DVC.} Pipeline \texttt{prepare → train → evaluate → export\_onnx} ;
données et poids versionnés par DVC, jamais par Git. \texttt{dvc repro} rejoue
le pipeline, \texttt{dvc metrics show} affiche les métriques.

\paragraph{MLflow.} Expérience \texttt{deepfake-adaptive-fusion} ; runs
comparables, artefacts, registry avec alias \texttt{staging}/\texttt{production}.
Le backend privilégie le modèle « Production » et retombe sur le fichier local.

\paragraph{API et UI.} FastAPI expose \texttt{/health},
\texttt{/model/metadata}, \texttt{/predict}, \texttt{/predict/visualize} avec
validation stricte (400/413/415/422/503), logs JSON et CORS. L'UI React propose
drag \& drop, galerie, webcam, barres animées, poids de fusion, Grad-CAM et
latences client/serveur.

\paragraph{Docker et CI.} Trois services (backend, frontend nginx, MLflow) via
\texttt{docker compose up --build}. Deux workflows GitHub Actions : tests
(black, flake8, pytest backend $\ge 80\,\%$, vitest, build) et build Docker avec
smoke test \texttt{/health} + \texttt{/predict}.

% -----------------------------------------------------------------------------
\section{Limites et perspectives}
\begin{itemize}[nosep]
  \item \textbf{Généralisation} : chute sur des générateurs inconnus (diffusion,
        StyleGAN3) et les vidéos fortement compressées.
  \item \textbf{Biais} : dataset majoritairement StyleGAN, qualité web,
        déséquilibre démographique probable.
  \item \textbf{Raccourcis} : artefacts de compression/générateur plutôt que
        sémantique du visage ; Grad-CAM n'est pas causal.
  \item \textbf{Seuil} : 0.5 par défaut, à recalibrer selon le coût métier.
  \item \textbf{Temps réel} : la webcam dépend de la lumière et de la qualité du
        flux ; pas de garantie forensique.
\end{itemize}
Perspectives : entraînement sur davantage de générateurs, apprentissage
contrastif, calibration par sous-groupe, détection vidéo temporelle, suivi de
dérive en production.

% -----------------------------------------------------------------------------
\section{Conclusion}
Le projet transforme un notebook de recherche en une chaîne MLOps complète :
données versionnées et anti-fuite, expériences tracées, modèle exporté et
quantifié, API testée, interface interactive, déploiement Docker et CI/CD. Les
résultats sont présentés sans embellissement, avec les limites de généralisation
et les risques d'usage. La reproductibilité (DVC, seeds, artefacts) et
l'honnêteté des conclusions sont les deux critères directeurs.

% -----------------------------------------------------------------------------
\appendix
\section{Annexes}
\subsection{Commandes de reproduction}
\begin{lstlisting}
git clone <repo> && cd <repo>
dvc pull
docker compose up --build
# tests
pytest backend -q --cov=backend/app --cov-fail-under=80
pytest tests -q -m "not slow"
# expériences
python scripts/run_experiments.py
python scripts/export_report_tables.py
\end{lstlisting}

\subsection{Captures}
Les captures MLflow, Swagger, frontend et CI/CD sont à déposer dans
\texttt{docs/} après exécution réelle, puis incluses ici via
\texttt{\textbackslash includegraphics}. Aucune capture n'est simulée.

\bibliographystyle{plain}
\bibliography{refs}

\end{document}
